% Informe del Laboratorio 1 (PSI 2-2026). Maximo 10 paginas en total.
%
% Cada seccion vive en su propio archivo en secciones/ y la escribe SOLO su responsable
% (ver el comentario al inicio de cada archivo). Asi dos personas nunca editan el mismo archivo
% y no hay conflictos de git. Este archivo principal no se toca para escribir una seccion.
%
% Las figuras NO se copian aqui: se toman de ../out/, que genera `psi-lab1`.
% Para compilar (desde la raiz del repositorio):
%     psi-lab1                              # genera out/ con todas las figuras y tablas
%     cd informe && latexmk -pdf lab1.tex   # genera informe/lab1.pdf
\documentclass[10pt]{article}
\usepackage[utf8]{inputenc}
\usepackage[T1]{fontenc}
\usepackage[spanish, es-noshorthands]{babel}
\usepackage[a4paper, margin=2cm]{geometry}
\usepackage{amsmath, amssymb}
\usepackage{graphicx}
\usepackage{booktabs}
\usepackage{float}
\renewcommand{\topfraction}{0.9}\renewcommand{\bottomfraction}{0.8}\renewcommand{\textfraction}{0.07}\renewcommand{\floatpagefraction}{0.85}\setcounter{topnumber}{3}\setcounter{totalnumber}{4}
\usepackage[section]{placeins}
\usepackage[hidelinks]{hyperref}

\graphicspath{{../out/}}  % figuras generadas por psi-lab1

% \figura{archivo.png}{ancho como fraccion de \linewidth}{leyenda}{etiqueta}
\newcommand{\figura}[4]{%
  \begin{figure}[htbp]
    \centering
    \includegraphics[width=#2\linewidth]{#1}
    \caption{#3}
    \label{#4}
  \end{figure}}

\title{Laboratorio 1: del muestreo al filtrado\\
       \large Procesamiento de Señales e Imágenes 2-2026 --- Caso C: imágenes}
\author{Dante Tureo Tapia \and Fernando Sarmiento Cortez \and Michel Vargas Carrasco \and Gustavo Vásquez Ortiz}
\date{Profesor: Miguel Cárcamo \\ 26 de noviembre de 2026}

\begin{document}
\maketitle

% Seccion: Objetivo e introducción
% Responsable: R1. Extension: 0,5 pagina(s).
% Solo R1 edita este archivo.
%
% Que debe contener (segun el enunciado y la pauta):
%   - Objetivo del laboratorio y del caso C.
%   - Organización del repositorio y cómo reproducir (docker compose run --rm lab1).
%   - Funciones auxiliares de bibliotecas usadas en el caso aplicado (np.fft, skimage: lectura y SSIM).

\section{Objetivo e introducción}
\label{sec:introduccion}

El laboratorio tiene dos objetivos. El primero es implementar desde cero, usando solo NumPy, las
herramientas de las Unidades 1--2 (muestreo, DFT, convolución y filtros FIR por ventana) y validarlas con
señales sintéticas cuyo resultado se conoce de antemano por la teoría. El segundo, el caso C, es usarlas
para estudiar Fourier en 2D: la DFT 2D y el papel de su magnitud y su fase, la eliminación de interferencia
periódica con filtros notch (con PSNR y SSIM contra la imagen original, y sobre fotos reales sin referencia)
y el aliasing espacial al reducir una imagen sin filtrarla antes.

\paragraph{Organización y reproducción.}
{\emergencystretch=3em  % evita que un nombre largo en \texttt se salga del margen
El paquete \texttt{psi\_lab.\allowbreak lab1} tiene un módulo por herramienta (\texttt{muestreo}, \texttt{dft},
\texttt{convolucion}, \texttt{filtros}, \texttt{ventanas}, \texttt{sinteticos}), un subpaquete
\texttt{caso\_imagen} para el caso C y un archivo por figura en \texttt{figuras/}. Cada figura y tabla de
este informe sale del código: \texttt{docker compose run --rm lab1} regenera todas en \texttt{out/}, y
descarga las dos fotos reales del caso C desde un \emph{release} del repositorio (no se versionan datos).
\texttt{docker compose run --rm tests} corre las pruebas unitarias (propiedades teóricas, no valores
copiados) y de integración; el CI ejecuta además el análisis estático y las pruebas en Python 3.11 a 3.13.\par}

\paragraph{Funciones de bibliotecas usadas.}
En la parte común solo se usa NumPy para el cálculo; \texttt{np.fft}, \texttt{np.convolve} y
\texttt{scipy} aparecen solo en las pruebas, como referencia (y \texttt{np.fft.fft} en la comparación
de tiempos de la DFT). En el caso C usamos como auxiliares \texttt{np.fft} para la DFT 2D y
\texttt{skimage} para cargar imágenes y calcular SSIM. La detección de picos, los notch, el PSNR y
el pasabajos gaussiano son propios.
   % R1  (0,5 pag.)
% Seccion: Muestreo y aliasing
% Responsable: R1. Extension: 1,0 pagina(s).
% Solo R1 edita este archivo.
%
% Que debe contener (segun el enunciado y la pauta):
%   - Deducción y justificación de la expresión de frecuencia_alias.
%   - Figura muestreo_alias.png y tabla muestreo_alias.csv: frecuencia medida vs. predicha.

\section{Muestreo y aliasing}
\label{sec:muestreo}

\paragraph{Frecuencia aparente.}
Una senoidal de frecuencia $f$ muestreada a $f_s$ da $x[n]=\cos(2\pi f n/f_s)$. Las muestras no cambian
(i)~al sumar a $f$ un múltiplo de $f_s$, porque $\cos(2\pi (f+kf_s)n/f_s)=\cos(2\pi fn/f_s+2\pi kn)$ y $kn$
es entero, ni (ii)~al cambiar su signo, porque el coseno es par. Todas las frecuencias $\pm f+kf_s$
producen entonces exactamente las mismas muestras y no se pueden distinguir. Con $k^\ast=\operatorname{round}(f/f_s)$,
la diferencia $f-k^\ast f_s$ cae en $[-f_s/2,\,f_s/2]$ y, por paridad, el representante del conjunto en
$[0,\,f_s/2]$ es
\begin{equation}
  f_a \;=\; \bigl|\,f - f_s\,\operatorname{round}(f/f_s)\,\bigr| ,
  \label{eq:alias}
\end{equation}
que es lo que calcula \texttt{frecuencia\_alias}. Si $f<f_s/2$ entonces $k^\ast=0$ y $f_a=f$ (sin aliasing,
criterio de Nyquist); si no, la frecuencia se refleja alrededor de los múltiplos de $f_s/2$.

\paragraph{Verificación.}
Se muestrearon senoidales de $f\in\{3,7,13,23\}$~Hz con $f_s=10$~Hz durante 10~s y se midió la frecuencia
aparente como el máximo del espectro (\texttt{espectro}, ventana rectangular, resolución
$\Delta f = 1/T = 0{,}1$~Hz). La ecuación~\eqref{eq:alias} predice 3~Hz en los cuatro casos, ya que
$7=10-3$, $13=10+3$ y $23=2\cdot10+3$; la medición coincide (tabla~\ref{tab:muestreo}). El error exacto
de cero se debe a que cada tono cae justo en un bin ($f$ y $T$ enteros); en general el error de la medición
está acotado por la resolución $\Delta f$.

% Valores de out/tablas/muestreo_alias.csv (generada por figuras/fig_muestreo.py).
\begin{table}[H]
  \centering
  \caption{Frecuencia aparente de senoidales muestreadas a $f_s=10$~Hz: predicha por
           \eqref{eq:alias} y medida con el espectro.}
  \label{tab:muestreo}
  \begin{tabular}{rrrr}
    \toprule
    $f$ [Hz] & $f_a$ predicha [Hz] & $f_a$ medida [Hz] & error [Hz] \\
    \midrule
    3  & 3,0 & 3,0 & 0,0 \\
    7  & 3,0 & 3,0 & 0,0 \\
    13 & 3,0 & 3,0 & 0,0 \\
    23 & 3,0 & 3,0 & 0,0 \\
    \bottomrule
  \end{tabular}
\end{table}

\figura{muestreo_alias.png}{0.72}{Las senoidales de 3~Hz y 13~Hz (curvas, muestreadas muy fino) pasan por
  exactamente las mismas muestras a $f_s=10$~Hz, pues $13=3+f_s$.}{fig:muestreo}

La figura~\ref{fig:muestreo} lo muestra para 3~Hz y 13~Hz: a $f_s=10$~Hz las muestras son idénticas, y por lo
tanto ningún procesamiento posterior puede recuperar cuál de las dos señales había. Las pruebas verifican
propiedades, no valores: $f_a\in[0,f_s/2]$, periodicidad en $f_s$, paridad e idempotencia de
\eqref{eq:alias}; que $f$ y $f_a$ generen las mismas muestras; y que, para 30 pares $(f,f_s)$ aleatorios,
el pico del espectro caiga en \texttt{frecuencia\_alias}.
       % R1  (1,0 pag.)
% Seccion: DFT, resolución y fuga espectral
% Responsable: R2. Extension: 1,5 pagina(s).
% Solo R2 edita este archivo.
%
% Que debe contener (segun el enunciado y la pauta):
%   - Ecuaciones de la DFT y de la normalización del espectro de un lado.
%   - 2.2 Tiempos (dft_tiempos.png): pendientes en log-log vs. O(N^2) y O(N log N).
%   - 2.3 Resolución (dft_resolucion.png): desde qué T se separan 50 y 52 Hz; Delta f = 1/T.
%   - 2.4 Relleno con ceros (dft_relleno.png): mejora la apariencia, no la resolución.
%   - 2.5 Fuga (dft_fuga.png, dft_fuga.csv): rectangular vs. Hann en dB con un tono de -60 dB.

\section{DFT, resolución y fuga espectral}
\label{sec:dft}

\textit{(Por escribir.)}
            % R2  (1,5 pag.)
% Seccion: Convolución directa y por FFT
% Responsable: R4. Extension: 1,0 pagina(s).
% Solo R4 edita este archivo.
%
% Que debe contener (segun el enunciado y la pauta):
%   - Ecuaciones de la convolución lineal y del teorema de la convolución.
%   - Tiempos de la directa vs. FFT según el largo del filtro (convolucion_tiempos.png).
%   - Convolución circular sin relleno de ceros (convolucion_circular.png).

\section{Convolución directa y por FFT}
\label{sec:convolucion}

La convolución lineal de $x$ (largo $N$) con $h$ (largo $M$) es
\begin{equation}
  y[n] = \sum_{k} x[k]\, h[n-k], \qquad n = 0, \dots, N+M-2,
  \label{eq:conv}
\end{equation}
y entrega $N+M-1$ muestras. \texttt{convolucion\_directa} la calcula con dos ciclos, y el ciclo de
adentro recorre $h$ para que el costo sea $O(NM)$.

\texttt{convolucion\_fft} usa el teorema de la convolución, que dice que convolucionar en el tiempo
es lo mismo que multiplicar en frecuencia. Si $\tilde{x}$ y $\tilde{h}$ son $x$ y $h$ completadas
con ceros hasta largo $L = N+M-1$,
\begin{equation}
  y = \mathrm{IDFT}\big(\mathrm{DFT}(\tilde{x}) \cdot \mathrm{DFT}(\tilde{h})\big),
  \label{eq:teorema-conv}
\end{equation}
donde el producto es punto a punto y usamos la \texttt{dft} y la \texttt{idft} de la
Sección~\ref{sec:dft}. Al final nos quedamos con la parte real, porque la parte imaginaria que
queda es casi cero ($10^{-16}$) y sale solo por el redondeo. En las pruebas comparamos las dos
versiones con \texttt{np.convolve} y entre ellas, y revisamos que un impulso en la posición $d$
mueve la señal $d$ muestras y que $x*h = h*x$.

\paragraph{Por qué rellenar con ceros.} La DFT de $L$ puntos supone que la señal se repite cada $L$
muestras, por eso al multiplicar DFTs sale la convolución \emph{circular}. Si no se rellena, las
últimas $M-1$ muestras no caben y se suman al comienzo. En la Figura~\ref{fig:conv}a usamos un
pulso cerca del final de $x$ ($N=40$) y una exponencial de $M=15$ coeficientes para que la cola
sea grande. Se ve que la cola de la lineal (muestras 40 a 53) aparece igual en las muestras 0 a 13
de la circular, donde la lineal vale cero.

\paragraph{Tiempos.} Dejamos $N=500$ fijo y probamos $M$ entre 4 y 512. Cada caso se midió 3 veces
y usamos el menor tiempo, para que no influyeran otros programas del computador. En escala log-log
la directa es una recta de pendiente 1, o sea que su tiempo crece igual que $M$, como dice $O(NM)$
(Figura~\ref{fig:conv}b). La versión por DFT casi no cambia, porque su costo depende de
$L = N+M-1$ y ese largo cambia poco mientras $M$ sea mucho menor que $N$. En nuestra prueba las
curvas se cruzan cerca de $M = 128$. Hay que tener cuidado con ese número, porque nuestra
\texttt{dft} es $O(L^2)$ pero usa operaciones rápidas de NumPy, mientras que la directa usa ciclos
de Python, que son más lentos. Por eso lo que conviene comparar son las pendientes. Con una FFT de
verdad, que es $O(L\log L)$, el cruce pasaría con filtros mucho más cortos.

\begin{figure}[H]
  \centering
  \begin{minipage}{0.55\linewidth}\centering
    \includegraphics[width=\linewidth]{convolucion_circular.png}\\(a)
  \end{minipage}\hfill
  \begin{minipage}{0.42\linewidth}\centering
    \includegraphics[width=\linewidth]{convolucion_tiempos.png}\\(b)
  \end{minipage}
  \caption{(a) Convolución lineal y circular sin relleno. La franja roja es la cola que se pasó
  al comienzo. (b) Tiempo de las dos versiones según el largo de $h$.}
  \label{fig:conv}
\end{figure}

Para el caso C agregamos \texttt{convolucion\_separable\_2d}, que pasa un filtro por cada fila y
otro por cada columna de la imagen. El resultado es el mismo que usar el filtro 2D
$h_\text{cols}\,h_\text{filas}^{T}$ (lo comparamos con \texttt{scipy.signal.convolve2d}), pero se
hacen $2M$ operaciones por píxel en vez de $M^2$. La usamos en el prefiltro de la
Sección~\ref{sec:aliasing}.
    % R4  (1,0 pag.)
% Seccion: Filtros FIR y su caracterización
% Responsable: R3. Extension: 1,5 pagina(s).
% Solo R3 edita este archivo.
%
% Que debe contener (segun el enunciado y la pauta):
%   - Diseño por ventanas; pasaaltos, pasabanda y rechazabanda por inversión espectral.
%   - Tabla de ancho de transición y atenuación según n_coef y ventana.
%   - Chirp filtrado (filtros_chirp.png).

\section{Filtros FIR y su caracterización}
\label{sec:filtros}

\textit{(Por escribir.)}
        % R3  (1,5 pag.)
% Seccion: Caso C: Fourier 2D y ruido periódico
% Responsable: R1 (DFT 2D, fase) y R2 (ruido periódico, detección). Extension: 1,5 pagina(s).
% Solo R1 (DFT 2D, fase) y R2 (ruido periódico, detección) edita este archivo.
% R1 escribe la subseccion "DFT 2D y fase"; R2 escribe "Ruido periódico y detección de picos".

\section{Caso C: Fourier 2D y ruido periódico}
\label{sec:fourier2d}

\subsection{DFT 2D y fase}
\label{sec:dft2}

\textit{(Por escribir, R1.)}

\subsection{Ruido periódico y detección de picos}
\label{sec:ruido}

Modelamos la interferencia periódica como un coseno 2D sumado a la imagen,
\begin{equation}
  r(x,y) = A\cos\!\left(2\pi\left(\frac{k_u\,x}{M} + \frac{k_v\,y}{N}\right)\right),
  \label{eq:coseno2d}
\end{equation}
con $x$ la columna, $y$ la fila e imagen de $N\times M$ píxeles. Elegimos $k_u$ y $k_v$ enteros,
es decir, un número exacto de ciclos en toda la imagen. Así la frecuencia cae justo en un bin de
la DFT 2D y aparecen exactamente dos picos, en $(k_v, k_u)$ y en su conjugado
$(N-k_v,\,M-k_u)$, cada uno de altura $A\,NM/2$. Si los ciclos no fueran enteros habría fuga, como en la
Sección~\ref{sec:dft}. Una prueba unitaria verifica que el coseno 2D da exactamente esos dos
picos.

Para encontrar los picos de forma automática (\texttt{detectar\_picos}), tomamos $|E|$, ponemos en
cero un disco de radio $r_{\text{DC}}$ alrededor del DC, nos quedamos con los bins que superan una
fracción \texttt{umbral} del máximo que queda, y agregamos siempre el conjugado de cada pico.
Esto último es necesario: si el notch borrara solo un pico del par, la imagen filtrada dejaría
de ser real.

El disco alrededor del DC es lo que hace funcionar esta versión simple. Una imagen natural tiene
casi toda su energía en frecuencias bajas, y sin el disco esos bins le ganan a cualquier
interferencia débil. Para elegir $r_{\text{DC}}$ probamos con la imagen \texttt{camera} y dos
interferencias de amplitud decreciente (Tabla~\ref{tab:sensibilidad}). Con $r_{\text{DC}}=10$
el detector encuentra los 4 picos sin ningún falso hasta una amplitud de 0{,}05; con
$r_{\text{DC}}=4$ ya aparecen falsos con 0{,}1. Bajo 0{,}02 ninguno de los dos funciona bien, y
ese es el límite del método. El costo es no detectar patrones de menos de 10 ciclos en la imagen,
que de todas formas se confunden con el contenido. Usamos \texttt{umbral} $=0{,}3$.

\begin{table}[H]
  \centering
  \small
  \begin{tabular}{ccccc}
    \toprule
    Amplitud & \multicolumn{2}{c}{$r_{\text{DC}}=4$} & \multicolumn{2}{c}{$r_{\text{DC}}=10$} \\
             & aciertos & falsos & aciertos & falsos \\
    \midrule
    0{,}2  & 4/4 & 0   & 4/4 & 0   \\
    0{,}1  & 4/4 & 4   & 4/4 & 0   \\
    0{,}05 & 4/4 & 50  & 4/4 & 0   \\
    0{,}02 & 4/4 & 106 & 4/4 & 140 \\
    0{,}01 & 0/4 & 106 & 4/4 & 634 \\
    \bottomrule
  \end{tabular}
  \caption{Picos detectados según la amplitud de la interferencia y el radio excluido alrededor
    del DC (\texttt{umbral} $=0{,}3$; tabla \texttt{ruido\_sensibilidad.csv}).}
  \label{tab:sensibilidad}
\end{table}

\figura{ruido_deteccion.png}{0.75}{Imagen \texttt{camera} limpia, con dos interferencias y su
  espectro centrado. Los círculos rojos son los 4 picos detectados (dos pares conjugados).}%
  {fig:ruido-deteccion}

En el espectro (Figura~\ref{fig:ruido-deteccion}) la interferencia aparece como puntos aislados,
mientras que el contenido se reparte en la mancha central y en líneas, que vienen de los bordes
rectos de la foto y de la cruz que forman los bordes de la imagen (la fuga de la ventana
rectangular, ahora en 2D). Por eso se puede borrar la interferencia sin tocar el resto.      % R1 + R2 (1,5 pag.)
% Seccion: Caso C: notch, métricas y caso real
% Responsable: R3 (notch, métricas) y R2 (caso real). Extension: 1,5 pagina(s).
% Solo R3 (notch, métricas) y R2 (caso real) edita este archivo.
%
% Que debe contener (segun el enunciado y la pauta):
%   - R3: notch ideal vs. gaussiano vs. pasabajos; PSNR y SSIM en varias condiciones.
%   - R2 + R3: caso real sin referencia (caso_real_diario.png, caso_real_pantalla.png, caso_real.csv).
%   - Cómo se evalúa sin imagen limpia; limitaciones (trama modulada, textura de la cámara).

\section{Caso C: notch, métricas y caso real}
\label{sec:notch}

\textit{(Por escribir.)}
          % R3  (1,5 pag.)
% Seccion: Caso C: aliasing espacial
% Responsable: R4. Extension: 1,0 pagina(s).
% Solo R4 edita este archivo.
%
% Que debe contener (segun el enunciado y la pauta):
%   - Zone plate y diezmado sin y con prefiltro; anillos fantasma predichos con frecuencia_alias.
%   - Justificación de la frecuencia de corte del prefiltro (fc ~ 0,5/k).

\section{Caso C: aliasing espacial}
\label{sec:aliasing}

\textit{(Por escribir.)}
       % R4  (1,0 pag.)
% Seccion: Discusión crítica y conclusiones
% Responsable: R4. Extension: 0,5 pagina(s).
% Solo R4 edita este archivo.
%
% Que debe contener (segun el enunciado y la pauta):
%   - Límites del método, casos de falla y alternativas.
%   - Conclusiones generales.

\section{Discusión crítica y conclusiones}
\label{sec:conclusiones}

\textit{(Por escribir.)}
   % R4  (0,5 pag.)

\end{document}
